\section{Análisis de errores}

El análisis se hizo a partir de las 5\,015 predicciones del holdout y no de una
selección visual de casos convenientes. La matriz de conteos permite seguir
cada error desde su clase real hasta la predicción. También se separaron los
errores con mayor confianza para comprobar si el problema detectado en la
Etapa 1 seguía presente.

\subsection{Deficiencias de nitrógeno, fósforo y potasio}

Las tres deficiencias reunieron 61 errores y 51 de ellos (83.6\%) terminaron en
otra clase del mismo bloque N/P/K. Nitrógeno tuvo 127 ejemplos: ocho se
predijeron como potasio y dos como fósforo. Fósforo tuvo 141: siete pasaron a
potasio y uno a nitrógeno. Potasio tuvo 93 y produjo el flujo más grande: 28 a
nitrógeno y cinco a fósforo. La confusión no desapareció al ampliar el corpus;
se volvió más asimétrica. El modelo suele absorber potasio dentro de nitrógeno,
pero también recibe falsos positivos de las otras deficiencias.

\subsection{Potasio: el cuello de botella anterior}

Solo 57 de las 93 imágenes de potasio fueron correctas. Su precision fue 0.7403,
recall 0.6129 y F1 0.6706. De los 36 errores, 28 terminaron en nitrógeno
(77.8\%), cinco en fósforo y uno respectivamente en gusano cogollero, sana y
necrosis letal. El problema principal no es una dispersión aleatoria por nueve
clases; es una frontera concreta entre amarillamientos nutricionales.

Frente al F1 histórico de 0.62 existe una mejora orientativa de 0.0506 y el
soporte de evaluación aumentó de 40 a 93. Sin embargo, el recall sigue dejando
sin reconocer casi cuatro de cada diez casos y en validación cruzada osciló con
desviación 0.0496. Potasio mejoró, pero continúa siendo el cuello de botella.

\subsection{Mancha gris y tizón norteño}

Mancha gris acertó 269 de 290 imágenes. Trece de sus 21 errores fueron hacia
tizón norteño. En el sentido inverso, 21 de 77 errores de tizón terminaron en
mancha gris; otros 25 se clasificaron como sanos. La relación visual detectada
en la Etapa 1 permanece, aunque sus F1 finales de 0.9150 y 0.9446 están muy por
encima del bloque nutricional. Aquí el modelo falla en una minoría de lesiones
parecidas; en potasio falla en la estructura central de la clase.

\subsection{Errores de alta confianza}

Siete de los 232 errores superaron 0.90 de confianza. El mayor fue una imagen
de fósforo predicha como potasio con 0.9779. Tres hojas enfermas ---dos con
gusano cogollero y una con tizón--- se mostraron como sanas con confianza entre
0.9296 y 0.9571; otra hoja de tizón recibió sana con 0.9145. También hubo roya
de campo como mancha gris (0.9178) y mancha gris de laboratorio como tizón
(0.9143).

Los casos más confiados cruzan precisamente dos límites de riesgo: confundir
insumos dentro de N/P/K y declarar sana una hoja afectada. El ensemble reduce
errores globales, pero no elimina la sobreconfianza. La interfaz debe mostrar
la segunda posibilidad, permitir rechazo y evitar convertir el porcentaje en
una orden de manejo.

La confianza máxima utilizada aquí procede de probabilidades one-vs-rest
normalizadas. Sirve para ordenar los casos, pero no está calibrada. Un error de
0.95 no demuestra que el sistema estuviera estadísticamente 95\% seguro; sí
demuestra que la interfaz habría mostrado una alternativa dominante. Esos
casos deben alimentar pruebas OOD, revisión de etiqueta y futuras corridas de
calibración.
